\documentclass[journal]{IEEEtran}

\usepackage[utf8]{inputenc}
\usepackage[T1]{fontenc}
\usepackage{amsmath,amssymb,amsthm}
\usepackage{booktabs}
\usepackage{graphicx}
\usepackage{array}
\usepackage{float}
\usepackage{caption}
\usepackage[colorlinks=true,linkcolor=blue,citecolor=blue,urlcolor=blue]{hyperref}
\usepackage{verbatim}
\usepackage[ruled,linesnumbered]{algorithm2e}
\usepackage{iftex}
\ifxetex
  \usepackage{xeCJK}
  \setCJKmainfont{FandolSong-Regular.otf}
\fi

\newtheorem{theorem}{Theorem}
\newtheorem{proposition}{Proposition}
\newtheorem{definition}{Definition}
\newtheorem{hypothesis}{Hypothesis}

\DeclareMathOperator*{\argmax}{arg\,max}

\title{Reversibility-Verified De-identification for Cloud-Local LLM Inference:\\
A Runtime-Sensitive-Set Release Chain with Layered Assurance (DR-SL)}
\author{Wen~Hu, Ya~Yu, and Xutong~Wang
\thanks{W. Hu is with the College of Electronic and Information Engineering, Nanjing University of Aeronautics and Astronautics, Nanjing, China, and with Jiangsu Yunhefeng Intelligent Technology Co., Ltd., Jiangsu, China (e-mail: huwen@nuaa.edu.cn). Y. Yu and X. Wang are with the College of Electronic and Information Engineering, Nanjing University of Aeronautics and Astronautics, Nanjing, China (e-mail: 15949116139@163.com; wsgxt@outlook.com). Corresponding author: Y. Yu.}
\thanks{A condensed conference version uses author order Y. Yu, W. Hu, X. Wang; the difference is disclosed in the submission declarations. Code and synthetic artifacts are public at \texttt{https://github.com/hurrywhen/dr-sl} and \texttt{https://gitee.com/jiangsu-yunhefeng-intelligent\_0/dr-sl}.}
\thanks{This revision uses fully synthetic data and co-author scoring of synthetic artifacts. Funding, ethics-review status, and venue-specific declarations must be finalized by the authors before submission.}}
\markboth{Manuscript under review}{Hu \MakeLowercase{\textit{et al.}}: Runtime-Sensitive-Set Release Chain for Cloud-Local LLM Inference}

\begin{document}
\maketitle

\begin{abstract}
Cloud-local LLM inference must keep sensitive user data on-device while exploiting cloud-grade reasoning. We study DR-SL, a release chain that combines runtime sensitive-set extraction, truthful generalization, local attack/utility checks, deterministic all-occurrence scrubbing, and human fallback. This revision corrects several overclaims in earlier versions: repeated probe copies are not independent Bernoulli trials; gold annotations are not a runtime sensitive set; reference-based online semantic judges are offline criteria, not local release inputs; and external re-testing is a controlled disclosure, not a pre-egress guarantee. We remove an invalid Fano threshold corollary, keep information-theoretic feasibility statements in mutual-information units, and implement the declared stability window. On a frozen independent-template holdout of 80 synthetic instances, a fully local deterministic runtime-S track releases 57.5\% of instances under a heuristic local gate, but external gold-literal auditing still finds residual literal exposure (mean gold hit rate 0.121) because runtime extraction recall is 0.858. A controlled-disclosure LLM audit on a 40-instance subset confirms residual hits (external max hit rate 0.130) and low semantic utility (0.363); the local utility surrogate correlates only weakly with the offline judge ($\rho=0.326$). Thus local release is not evidence of zero egress leakage. Oracle-S experiments are reported separately as upper-bound references. Human-study reanalysis shows overall utility correlation $\rho=0.839$ but substantially lower within-operating-point correlations, and the original combination-risk annotation criterion failed ($\kappa=-0.083$); a construct-aligned $\kappa=0.696$ is post hoc. Self-learning remains exploratory: its preregistered efficiency hypothesis was not supported. The contribution is therefore an auditable cloud-local release workflow and a clear account of its failure modes, not a certified zero-leakage or fully automatic system.

\medskip
\noindent\textbf{Keywords:} sensitive-information isolation; cloud-local inference; runtime sensitive set; release decision; truthful generalization; Pufferfish privacy; controlled-disclosure audit
\end{abstract}

\section{Introduction}

Cloud models offer stronger reasoning than on-device small models, but forwarding user requests exposes sensitive data. Cloud-local systems seek both capabilities under the constraint that cloud-bound requests should exclude personally identifiable information (PII)~\cite{hao2024hybrid,jin2025cecollm,li2025collab,siyan2025papillon,das2025security,lukas2023pii}. The difficulty is not only removing names or ID numbers: quasi-identifier combinations can identify a person, while deleting task-critical context can make the cloud response useless.

This paper revises DR-SL around a narrower, evidence-supported claim: \emph{an auditable local release chain can combine runtime sensitive-set extraction, truthful generalization, attack/utility checks, deterministic scrubbing, and human fallback, but its safety is conditional on extraction coverage, attacker class, and utility calibration}. We do not claim statistical zero leakage, fully automatic release, or that self-learning is the source of the final safe operating point.

\subsection{Corrections relative to earlier versions}

The revision makes the following changes:
\begin{itemize}
\item \textbf{Statistical unit.} Earlier experiments replicated the same question bundle to reach $\ge75$ probes and reused one enumeration response across replicated questions. Those observations are not independent. The revised gate is a heuristic observed-hit-rate gate; instance-cluster bootstrap replaces pooled binomial/Wilson inference for attack summaries.
\item \textbf{Runtime-S.} Earlier core experiments used gold annotations inside dehydration, strengthening, QA, and hard-line scrubbing. The revised runtime path receives only extractor output and a declared task-relevance profile. Gold is used only by external scoring after the artifact and local decision are frozen.
\item \textbf{Local vs online utility.} The earlier mixed-benchmark ``automatic release'' used an online reference-based judge that saw original $X$. That is offline oracle routing, not fully local release. The revised local gate uses a fully local utility surrogate; online judges are offline criteria only.
\item \textbf{Theory.} We delete the invalid Fano threshold-mapping corollary, avoid substituting QA scores or hit-rate thresholds for mutual-information budgets, and prove termination from finite budgets rather than bounded continuous improvement.
\item \textbf{External re-test.} Sending $D$ to an external attacker is itself a disclosure. We reclassify it as a controlled-disclosure audit, not a pre-egress safety boundary.
\item \textbf{Baselines and attribution.} We fix first-occurrence-only replacement in the random baseline, rename it PRISM-style random replacement rather than LDP, add an INTACT-inspired candidate-selection baseline, and report deterministic/generalization/loop ablations under the same runtime-S and gates.
\end{itemize}

\subsection{Contributions}

\begin{itemize}
\item A runtime-S release chain with explicit trust boundaries, stable local gates, deterministic all-occurrence scrubbing, residual reports, and human fallback.
\item A corrected theoretical account: Fano-type leakage witnesses at actual MAP success, a Pufferfish witness, an MI-unit feasibility inequality, and finite-budget termination.
\item A frozen independent-template holdout and fair runtime-S ablations, with oracle-S results separated from deployment-track results.
\item A restrained human-study reanalysis reporting within-group correlations, instance-cluster intervals, MAE, threshold confusion, and the failed original annotation criterion.
\end{itemize}

\section{Related Work and Positioning}

\paragraph{Cloud-local inference and routing.} PAPILLON formulates privacy-conscious delegation with a local proxy deciding what to forward~\cite{siyan2025papillon}. Hybrid and collaborative inference systems study capability/latency trade-offs~\cite{hao2024hybrid,jin2025cecollm,li2025collab}. DR-SL differs by making the release decision depend on a runtime sensitive set and explicit residual reporting, but it does not claim to solve routing generally.

\paragraph{INTACT is the closest methodological baseline.} INTACT~\cite{pilan2025truthful} uses instruction-tuned open-weight LLMs to generate truth-preserving generalization candidates at different abstraction levels, then selects the most informative candidate resistant to inference attacks. This is very close to DR-SL's generalization--attack--utility idea. We therefore do not claim that candidate generalization plus attack selection is novel. Our incremental scope is the cloud-local release chain: runtime-S construction, explicit/quasi-identifier/task-information joint gating, INFEASIBLE residual reports, human routing, and probe-format failure analysis. Because the official INTACT implementation was not run here, our \texttt{intact\_style} baseline is an algorithmic proxy, not a reproduction.

\paragraph{Placeholder and restoration systems.} Presidio and Rehydra provide entity detection/substitution and restoration workflows~\cite{presidio,rehydra}. PrivacyRestore removes privacy spans client-side and restores server-side via meta-vectors~\cite{zeng2025privacyrestore}. These systems motivate reversible handling but do not by themselves establish quasi-identifier release decisions.

\paragraph{Perturbation and skill distillation.} LDP-style text perturbation trades privacy for coherence~\cite{mai2023split,shi2022selective,duchi2013local,yue2021santext}. Our random same-type replacement baseline is PRISM-style, not a formal LDP mechanism. P2Skill distills reusable privacy skills into a frozen local SLM~\cite{ryu2026p2skill}; its refinement loop uses cloud reference/evaluator roles on original inputs, while deployment egress is matcher-gated. The earlier $1.69\times/3.66\times$ numbers come from v1; v2 (2026-10-13) changes evaluation emphasis and should be cited separately. We therefore avoid the claim that existing methods lack cross-instance reuse: P2Skill explicitly distills skills. The narrower gap we address is joint runtime-S release gating and residual routing for cloud-local inference.

\section{Threat Model and Runtime Sensitive Set}

Let $X$ be original material, $S$ a sensitive attribute set, $A$ attacker auxiliary knowledge, $T$ task-relevant information, and $D$ the released artifact. The local device is trusted; the cloud/link adversary observes egressed artifacts. A compromised local device voids the guarantees.

\subsection{Runtime-S versus oracle-S}

The revised runtime path constructs $S_{\mathrm{rt}}$ from $X$ using an extractor. Policy functions receive only $(X,\mathrm{domain},S_{\mathrm{rt}},\Pi)$, where $\Pi$ is a declared task-relevance profile. Gold annotations $G$ are read only by external scoring. We report two tracks:
\begin{itemize}
\item \textbf{oracle-S}: $S=G$ enters the policy. This is an upper-bound reference, not deployment evidence.
\item \textbf{runtime-S}: $S=S_{\mathrm{rt}}$ enters the policy. Missed gold attributes are outside the local gate and must be measured externally.
\end{itemize}

The deterministic runtime extractor used in the new experiments combines explicit patterns, public synthetic-domain lexicons, and generic labeled-span rules. It is scoped to the synthetic benchmark; it is not an open-world NER guarantee. The neural cascade extractor remains implemented for local-LLM deployment, but the local neural service was unavailable for this revision, so no neural local-LLM end-to-end rerun is claimed.

\subsection{Recall sensitivity, not an average upper bound}

For a fixed instance, let $F$ be found attributes and $G$ gold attributes. If probes cover $F$ and $\mathrm{hit}(s)$ equals the leakage indicator for $s\in F\cap G$, while missed attributes contribute at most 1, then
\[
L_{\mathrm{system}} \le \frac{|F|}{|G|} L_{\mathrm{measured}} + (1-\mathrm{Recall}).
\]
This is a per-instance algebraic sensitivity identity under those assumptions. Substituting \emph{average} recall and average over-extraction ratios does not yield a valid system-wide upper bound. We therefore report recall, over-extraction, missed-gold exposure, and external gold-literal hits separately.

\section{Release Chain and Local Gates}

\subsection{Policies}

All methods share the same runtime entities, task profile, local attack strategies, utility surrogate, and release rule:
\begin{itemize}
\item \texttt{hardline\_only}: typed placeholders for every runtime entity and every occurrence.
\item \texttt{generalization\_only}: truthful numeric/date bands for task-relevant values, placeholders otherwise.
\item \texttt{closed\_loop}: generalization followed by the stable local gate, final all-occurrence scrub, and gate recheck.
\item \texttt{intact\_style}: INTACT-inspired candidate selection among placeholder/generalized/loosened artifacts, choosing the highest local utility among candidates passing the local privacy gate.
\item \texttt{prism\_random}: consistent random same-type replacement for every occurrence; no DP guarantee is claimed.
\end{itemize}

\subsection{Local attack and utility}

The local privacy audit is deterministic and fully local. It checks verbatim, punctuation-normalized, and NFKC-normalized occurrences of runtime entity values. Strategy hit rates are not concatenated into a larger sample; the maximum strategy rate is the conservative local summary. The gate is
\[
L_{\mathrm{local}} \le \tau_L,\qquad \tau_L=0.05,
\]
where $L_{\mathrm{local}}$ is an observed point rate, not a confidence bound.

The local utility surrogate is reference-based but local: exact retention scores 1; a frozen truthful numeric/date band scores 0.5; deletion, placeholders, or wrong random values score 0 for each runtime task-relevant attribute. Instances without such attributes use non-sensitive context retention. The local utility gate is $U_{\mathrm{local}}\ge0.5$. This surrogate is not a human-equivalent semantic judge.

\subsection{Stability and termination}

The declared stability window is implemented: $N=3$ evaluations of the \emph{same frozen artifact}, all passing, with ranges of $L$ and $U$ at most $\delta=0.05$. A changed artifact resets the window. The returned artifact must itself have a complete stable window.

\begin{algorithm}[!t]
\caption{Revised local release loop (schematic).}
\KwIn{$X$, runtime entities $S$, task profile $\Pi$, thresholds $\tau_L,\tau_U$, budget $B=3|S|+5$, $N=3$, $\delta=0.05$}
\KwOut{PASS($D$) or INFEASIBLE(residual report)}
$D\leftarrow$ generalize/scrub according to policy\;
evaluate $D$; initialize window\;
\While{utility phase needs loosening and $t<B$}{
  propose loosen; increment $t$; accept only if diagnostic improves; reset window on artifact change\;
}
\While{not stable and $t<B$}{
  \If{current artifact passes and window incomplete}{
    re-evaluate same artifact; update window\;
  }
  choose unskipped attribute with positive local contribution; increment $t$\;
  propose single-attribute strengthening; accept only if utility gate holds; reset window on change\;
  otherwise skip attribute or record residual\;
}
apply final all-occurrence runtime scrub; recheck stable gates\;
\eIf{stable PASS}{\Return PASS($D$)}{\Return INFEASIBLE(residual report)}
\end{algorithm}

Termination follows from the finite round budget and finite skip set, not from boundedness of a continuous utility score.

\subsection{External audit is controlled disclosure}

After local decisions are frozen, external attackers may receive $D$. This is a first disclosure to the audit trust class. A failed audit can prevent subsequent processing but cannot undo the first exposure. Therefore external re-testing is not a pre-egress guarantee. A local strong attacker or trusted execution environment would be needed to avoid this disclosure; those are future deployment options.

\section{Theoretical Properties}

\subsection{Fano-type leakage witness}

\begin{theorem}[Fano-type lower bound at MAP success]
Let $S=(S_1,\dots,S_m)$ have finite domains $V_i$, and let $L_i^{*}$ be the optimal MAP decoding success probability for $S_i$ from $D$ under a specified joint distribution and prior. Then
\[
I(D;S) \ge H(S)-\sum_i \left[H_2(1-L_i^{*})+(1-L_i^{*})\log(|V_i|-1)\right].
\]
\end{theorem}

The bound certifies leakage, not safety. A local evaluator's failure to recover a value does not prove that no stronger attacker can. Empirical LLM template hit rates are not automatically MAP success probabilities under a defined random trial; we use them as heuristic attack coverage unless a sampling model is established.

\paragraph{Deleted threshold corollary.} The earlier claim that $L_i\le\tau_L$ implies $I(D;S)\ge H(S)-\sum_i g(1-\tau_L)$ is false. Counterexample: $S$ uniform on 100 values, $D$ independent of $S$, so $I(D;S)=0$ and optimal guessing success is $0.01\le0.05$, yet substituting $\tau_L=0.05$ gives a positive lower bound. The regression test \texttt{test\_fano\_threshold\_corollary\_counterexample} records this.

\subsection{Pufferfish witness}

For a single attribute with prior extremes $p_{\min},p_{\max}$ and MAP success $L^{*}$, any mechanism satisfying Pufferfish odds constraints must have
\[
\varepsilon \ge \ln\frac{L^{*}}{1-L^{*}}+\ln\frac{p_{\min}}{p_{\max}}.
\]
This is a leakage witness. It is nontrivial only when the attacker is strong; low hit rates give vacuous bounds.

\subsection{MI-unit feasibility inequality}

Assume $T$ is a function of $X$, and conditioned on $(S,A)$ the artifact satisfies $T\to X\to D$. Then
\[
I(D;T\mid A)\le I(D;S\mid A)+I(X;T\mid S,A).
\]
Consequently, for an MI privacy budget $\varepsilon$ and MI utility requirement $\upsilon$, the feasible region is empty if $\upsilon>\varepsilon+I(X;T\mid S,A)$. We do \emph{not} substitute QA scores or attack hit-rate thresholds for $\varepsilon$ or $\upsilon$. Low medical release rates are consistent with strong task--sensitive coupling qualitatively, but do not confirm the theorem quantitatively.

\subsection{Numerical tests are regression tests}

The proof-check suite performs exact-enumeration regression and counterexample search on synthetic mechanisms. It is not an independent proof of theorem correctness. The revised suite includes the deleted-corollary counterexample and finite-budget/stability tests.

\section{Experiments}

\subsection{Protocol and data separation}

The revised protocol is frozen in \texttt{revision\_protocol.md} before holdout evaluation. Earlier T1--T25/P1--P41 results are method-development evidence only. The new holdout contains 80 synthetic instances (20 per domain), seed 20261118, generated from an independent template family (\texttt{dataset/generate\_holdout.py}). The external audit subset is the first 10 holdout instances per domain after sorting by ID. No threshold, prompt, lexicon, or template is tuned after holdout results are inspected; process corrections are logged in \texttt{revision\_deviations.md}.

All values are fictional. MedQA-derived stems in the development set use MIT-licensed exam questions with concrete values replaced~\cite{jin2021medqa}. No real patient data are used.

\subsection{Metrics and statistics}

We report local observed hit rate, local utility surrogate, release rate, extraction recall/over-extraction, missed-gold exposure, and external gold-literal hit rate. Attack summaries use instance-cluster bootstrap (2000 resamples, seed 20261118). Release-rate Wilson intervals are used only at the independent-instance level. Permutation p-values use $(b+1)/(B+1)$; $p=0$ is forbidden. Zero hits are reported as zero \emph{observed} hits under specified attackers and grading, with coverage counts, not as a statistical guarantee.

\subsection{Runtime-S and oracle-S results}

% Auto-generated by sim/make_revision_tables.py (round 3). Do not edit by hand.

\begin{table}[!t]\centering\footnotesize
\caption{Holdout-1 closed loop stratified by local decision (Wilson 95\% intervals). The primary safety metric is the conditional residue rate given local PASS (FR$_{gold}$); whole-set averages mask the released subgroup.}
\label{tab:stratified}
\begin{tabular}{@{}lcccc@{}}\toprule
Stratum & $n$ & residue \% & 95\% CI & gold $L$ \\ \midrule
ALL & 80 & 68.8 & [.579,.778] & 0.121 \\
PASS & 46 & 78.3 & [.644,.877] & 0.147 \\
INFEASIBLE & 34 & 55.9 & [.395,.711] & 0.086 \\
\bottomrule\end{tabular}\end{table}

\begin{table}[!t]\centering\footnotesize
\caption{Holdout-1 closed loop by domain. Banking releases 20/20 with residue in every released instance; legal PASS is 10/10 with residue. All 20 banking instances enter the harmful utility fallback (mechanism analysis in text).}
\label{tab:domain}
\begin{tabular}{@{}lccccc@{}}\toprule
Domain & rel.\ \% & PASS $n$ & FR$_{gold}$ \% & 95\% CI & PASS gold $L$ \\ \midrule
banking & 100.0 & 20 & 100.0 & [.839,1.000] & 0.211 \\
legal & 50.0 & 10 & 100.0 & [.722,1.000] & 0.155 \\
medical & 25.0 & 5 & 0.0 & [.000,.434] & 0.000 \\
travel & 55.0 & 11 & 54.5 & [.280,.787] & 0.091 \\
\bottomrule\end{tabular}\end{table}

\begin{table}[!t]\centering\setlength{\tabcolsep}{2pt}\footnotesize
\caption{Holdout-1 method comparison under the same local gate (40-instance audited subset for external columns). HL=hardline-only, GEN=generalization-only, CL=closed loop, INT=INTACT-style proxy, PRISM=random replacement, DF=direct-forward reference. ext=external max hit rate; judge=offline artifact-judge utility; task=reference-answer preservation (primary judge). Paired rows: CL$-$GEN; the task-$U$ delta uses the round-3 rerun. Local-utility delta is negative on all three holdouts; release/leakage deltas are zero throughout.}
\label{tab:methods}
\begin{tabular}{@{}lcccccc@{}}\toprule
Method & rel.\ \% & $U_{loc}$ & gold $L$ & ext $L$ & judge $U$ & task $U$ \\ \midrule
HL & 51.2 & 0.512 & 0.121 & 0.130 & 0.425 & -- \\
GEN & 57.5 & 0.632 & 0.121 & 0.130 & 0.438 & 0.700 \\
CL & 57.5 & 0.544 & 0.121 & 0.130 & 0.362 & 0.725 \\
INT & 57.5 & 0.632 & 0.121 & 0.130 & 0.412 & -- \\
PRISM & 51.2 & 0.537 & 0.132 & 0.139 & 0.125 & -- \\
DF & -- & -- & -- & 1.000 & 1.000 & -- \\
\multicolumn{7}{@{}l@{}}{\emph{Paired $\Delta$ CL$-$GEN; release/gold-$L$ deltas 0; CIs in text}} \\
$\Delta$ h1 & 0.0 & -0.089 & 0.0 & 0.0 & -0.075 & 0.025 \\
$\Delta$ h2 & 0.0 & -0.017 & 0.0 & -- & -- & -- \\
$\Delta$ h3 & 0.0 & -0.101 & 0.0 & -- & -- & -- \\
\bottomrule\end{tabular}\end{table}

\begin{table}[!t]\centering\footnotesize
\caption{Vocabulary-shift stress and dose--response (closed loop, original extractor). overlap = realized fraction of lexicon-dependent gold values from the old vocabulary. holdout-2 is a constructive fully-disjoint lexicon ablation, not a deployment failure rate; the drift series holds templates constant and varies only overlap (40 instances/level).}
\label{tab:vocab}
\begin{tabular}{@{}lccccc@{}}\toprule
Condition & overlap & recall & rel.\ \% & FR$_{gold}$ \% & 95\% CI \\ \midrule
holdout-1 & 1.000 & 0.858 & 57.5 & 78.3 & [.644,.877] \\
holdout-2 & 0.000 & 0.526 & 93.8 & 100.0 & [.951,1.000] \\
drift p000 & 0.000 & 0.522 & 92.5 & 100.0 & [.906,1.000] \\
drift p025 & 0.163 & 0.582 & 80.0 & 100.0 & [.893,1.000] \\
drift p050 & 0.408 & 0.675 & 77.5 & 93.5 & [.793,.982] \\
drift p075 & 0.682 & 0.745 & 65.0 & 92.3 & [.759,.979] \\
drift p100 & 1.000 & 0.845 & 57.5 & 82.6 & [.629,.930] \\
\bottomrule\end{tabular}\end{table}

\begin{table}[!t]\centering\setlength{\tabcolsep}{3pt}\footnotesize
\caption{Fixed-extractor contrast. The fixed variant repairs clear defects only (money pattern wiring, ASCII email local part, allergy label/collision); no lexicon additions. Holdout-3 (natural 50/50 vocabulary mix, new templates) is confirmatory. CG = coverage gate (UNSCORABLE$\to$human when no task attribute extracted). FB = fallback \% all/benign/harmful (benign: no gold task attributes either). Fixing defects leaves no observed gold false release on shared-vocabulary h1 (0/46) but not under vocabulary shift. Task recall: h1 1.000, h2 0.750, h3 0.427$\to$0.792.}
\label{tab:fixcontrast}
\begin{tabular}{@{}lcccc@{}}\toprule
Config & recall & rel.\ \% & FR$_{gold}$ \% [95\% CI] & FB \% a/b/h \\ \midrule
h3 original & 0.542 & 55.0 & 100.0 [.920,1.000] & 30.0/0.0/30.0 \\
h3 fixed & 0.713 & 51.2 & 82.9 [.687,.915] & 5.0/0.0/5.0 \\
h3 orig.+CG & 0.542 & 25.0 & 100.0 [.839,1.000] & 30.0/0.0/30.0 \\
h3 fixed+CG & 0.713 & 46.2 & 81.1 [.658,.905] & 5.0/0.0/5.0 \\
\multicolumn{5}{@{}l@{}}{\emph{Diagnostic (defects discovered on these sets; not confirmatory)}} \\
h1 fixed & 1.000 & 57.5 & 0.0 [.000,.077] & 13.8/13.8/0.0 \\
h2 fixed & 0.668 & 87.5 & 100.0 [.948,1.000] & 18.8/12.5/6.2 \\
\bottomrule\end{tabular}\end{table}

\begin{table*}[!t]\centering\setlength{\tabcolsep}{2.5pt}\footnotesize
\caption{Cloud-assisted neural-extractor contrast (qwen3.8-max, frozen prompt; DIAGNOSTIC ONLY---extraction discloses $X$; single-run, instance-level intervals). Same gates/methods; det.=original deterministic extractor (decomposition n/a). FR split by residual gold type: gap=undeclared type (room\_no, marriage\_years, work\_years omitted on h1/drift; h3 fully declared); decl=within-taxonomy miss. h1/drift-p0 residue is entirely room\_no gap; only h3 evidences declared-type neural misses. h4 neural v2 uses the complete 43-type taxonomy (prompt v2, round 5) on NEW holdout-4 (template family \#4, natural mix). Coverage layers and bound scoping in text: exact-span recall is conservative under the prompt/taxonomy; FR is conditional on the incomplete taxonomy, with no general bound direction.}
\label{tab:neural}
\begin{tabular}{@{}lcccccc@{}}\toprule
Data & recall & task rec. & rel.\ \% & FR$_{gold}$ \% [95\% CI] & FR$_{gap}$ \% [95\% CI] & FR$_{decl}$ \% [95\% CI] \\ \midrule
h1 det. & 0.858 & 0.493 & 57.5 & 78.3 [.644,.877] & -- & -- \\
h1 neural & 0.965 & 1.000 & 38.8 & 16.1 [.071,.326] & 16.1 [.071,.326] & 0.0 [.000,.110] \\
h3 det. & 0.542 & 0.427 & 55.0 & 100.0 [.920,1.000] & -- & -- \\
h3 neural & 0.973 & 0.971 & 30.0 & 16.7 [.067,.359] & 0.0 [.000,.138] & 16.7 [.067,.359] \\
drift p0 det. & 0.522 & 0.229 & 92.5 & 100.0 [.906,1.000] & -- & -- \\
drift p0 neural & 0.965 & 1.000 & 37.5 & 13.3 [.037,.379] & 13.3 [.037,.379] & 0.0 [.000,.204] \\
h4 det. & 0.553 & 0.359 & 63.7 & 100.0 [.930,1.000] & -- & -- \\
h4 neural v2 & 0.986 & 0.994 & 40.0 & 3.1 [.006,.157] & 0.0 [.000,.107] & 3.1 [.006,.157] \\
\bottomrule\end{tabular}\end{table*}

\begin{table}[!t]\centering\footnotesize
\caption{Utility-gate fallback mechanism (closed loop, original extractor). Harmful fallback: gold task-relevant attributes exist but none were extracted, so local utility scores unmarked-context retention instead (context the extractor did not flag, which may contain missed sensitive values). FR$_{gold}$ among released harmful-fallback instances vs 0.375 among other released instances on h1: worse extraction makes release easier.}
\label{tab:mechanism}
\begin{tabular}{@{}lccccc@{}}\toprule
Data & attr rec. & task rec. & harmful FB \% & rel.\ via FB & FR$_{gold}$|FB \\ \midrule
h1 & 0.858 & 0.493 & 37.5 & 30/46 & 1.000 \\
h2 & 0.526 & 0.250 & 43.8 & 35/75 & 1.000 \\
h3 & 0.542 & 0.427 & 30.0 & 24/44 & 1.000 \\
\bottomrule\end{tabular}\end{table}

\begin{table}[!t]\centering\footnotesize
\caption{Utility measurement (round-3 corrected, round-4 direction fix). Calibration: local surrogate vs offline artifact judge, instance-cluster 95\% CIs, average tied ranks. Task $U$ = reference-answer preservation (agreement with the answer from full $X$; not absolute correctness), explicit tasks, hash-cached answers; primary judge online\_alt, robustness judge online\_fast, both reported. The surrogate is uncalibrated in absolute error and threshold decisions; rank correlation is moderate with wide intervals.}
\label{tab:utility}
\begin{tabular}{@{}lcc@{}}\toprule
Quantity & estimate & 95\% CI \\ \midrule
Spearman $\rho$ (tied ranks) & 0.450 & [0.242,0.634] \\
MAE & 0.352 & [0.265,0.448] \\
threshold agreement & 0.680 & [0.565,0.790] \\
CL confusion TP/FP/FN/TN & 17/6/5/12 & -- \\
task $U$: CL (ALL) & 0.725 & [0.613,0.838] \\
task $U$: CL (PASS / INFEAS) & 0.826 / 0.588 & -- \\
task $U$: GEN (ALL) & 0.700 & [0.575,0.825] \\
task $U$: GEN (PASS / INFEAS) & 0.826 / 0.529 & -- \\
paired $\Delta$ CL$-$GEN, primary & 0.025 & [0.0000,0.0625] \\
paired $\Delta$ CL$-$GEN, robust. & 0.025 & [-0.0375,0.0875] \\
score disagreement (prim./robust.) & 2/40 vs 8/40 & -- \\
judge exact agreement / MAE & 0.725 / 0.138 & -- \\
\bottomrule\end{tabular}\end{table}

\begin{table}[!t]\centering\footnotesize
\caption{Constructed two-type residue probe (55 holdout-1 instances with gold residue). Type a: an extracted value re-inserted (exact = verbatim; evasion = normalized-equivalent variant that exact replacement misses). Type b: never-extracted gold residue. Report coverage is measured among INFEASIBLE loop exits; PASS instances have no report by construction. Detectable residue is reported/routed; never-extracted residue is invisible.}
\label{tab:residual}
\begin{tabular}{@{}lcccc@{}}\toprule
Mode & type a detect & type a report & type a scrubbed & type b report \\ \midrule
exact & 1.000 & 0.000 & 1.000 & 0.000 \\
evasion & 1.000 & 0.655 & 0.273 & 0.000 \\
\bottomrule\end{tabular}\end{table}

\begin{table}[!t]\centering\footnotesize
\caption{Open-ended inference attack (no gold question types; exploratory FR$_{open}$, still literal-graded). DF is a reference under this protocol, not a theoretical attacker ceiling. Released artifacts remain highly recoverable even without question-type hints.}
\label{tab:openattack}
\begin{tabular}{@{}lcccc@{}}\toprule
Method & FR$_{open}$ \% & 95\% CI & FR$_{open}$|PASS \% & mean $L$ \\ \midrule
CL & 72.5 & [.572,.839] & 82.6 & 0.135 \\
GEN & 67.5 & [.520,.799] & 78.3 & 0.121 \\
DF & 100.0 & [.912,1.000] & -- & 0.996 \\
\bottomrule\end{tabular}\end{table}


Table~\ref{tab:runtime_oracle} separates oracle-S from runtime-S. On the fully coupled development test-160, oracle-S local utility is too low for automatic release under the local surrogate. On mixed-120, oracle-S release is high, but runtime-S release must be interpreted together with external gold-literal exposure. On the frozen holdout, runtime extraction recall is 0.858 and missed gold attributes produce residual literal exposure even when the local gate passes. This is the central negative result: \emph{local release under runtime-S is not evidence of zero egress leakage}.

Table~\ref{tab:holdout_methods} compares fair runtime-S methods. The closed loop and generalization-only artifact are often identical because the deterministic generalizer already handles all occurrences; the loop's incremental value in this deterministic track is stable gating and residual routing, not a new safe artifact. The INTACT-inspired selector matches the generalized point under the same local gate, reinforcing that candidate selection itself is not the novel contribution. PRISM-style random replacement is not labelled LDP and does not dominate the truthful generalization point.

The controlled-disclosure audit (Table~\ref{tab:audit_coverage}) used two online attackers and two strategies on 40 frozen holdout instances, with 640 model calls and 3888 unique questions; it observed 519 hits, not zero. For the closed-loop artifact, external max hit rate is 0.130 [0.103, 0.159] and offline semantic-judge utility is 0.363 [0.250, 0.475]. The local utility surrogate is poorly calibrated against this offline judge: overall Spearman $\rho=0.326$, MAE 0.358, threshold agreement 68.1\%; within-method correlations range from $-0.071$ (random replacement) to 0.535 (INTACT-style). Table~\ref{tab:threshold_confusion} shows 31 artifacts released locally but judged below threshold offline. These results reinforce that the local gate is a heuristic routing signal, not a validated safety or utility certificate.

\subsection{Historical exploratory results}

Earlier oracle-S experiments reported leakage reduction $0.457\to0.304$ on train-400 and PASS rates 37.5\%/43.8\%/5.0\% under different probe formats. Those numbers used replicated probe bundles and gold-S policies; they are not confirmatory under the revised gate. The earlier mixed-benchmark 67.5\% was an offline oracle semantic-routing rate because the online judge saw original $X$. We retain these results only as exploratory development history.

\subsection{Human studies}

Utility study: two raters scored 160 blinded rows (40 instances $\times$ four operating points). Inter-rater Krippendorff $\alpha=0.780$. Overall human--automatic Spearman $\rho=0.839$ with instance-cluster 95\% CI [0.774, 0.892]; MAE 0.133 [0.100, 0.169]. Within-operating-point correlations are lower: generalized 0.661 [0.442, 0.830], placeholder 0.784 [0.587, 0.937], random replacement 0.626 [0.349, 0.812]; direct has no variance. The four-point rank correlation is descriptive only. Threshold confusion around 0.5 is reported in the analysis package. Thus the semantic judge has some overall ranking validity but is not a fully calibrated local release threshold.

Annotation study: type-level recall against automatic gold is 1.000/0.987 and precision 0.959/0.991 for the two raters. The original combination-risk criterion failed: human--automatic $\kappa=-0.083$. A construct-aligned gold (combos or any explicit identifier) gives $\kappa=0.696$, but this is post hoc and does not repair the preregistered failure. Human recall against automatic gold also does not prove the automatic gold is correct.

\subsection{Self-learning is exploratory}

The preregistered loop-efficiency hypothesis H1 was not supported: distilled rules changed validation rounds from 2.82 to 2.90. The corrected extractor-targeted hypothesis H1$'$ showed a small directional recall gain ($+1.5$ pp) but did not reach significance ($p=0.059$). We therefore do not claim cross-instance self-learning as the source of the final safe point. The supported components are deterministic generalization, all-occurrence scrubbing, stable local gating, and residual routing.

\section{Limitations}

\begin{itemize}
\item \textbf{Extraction coverage.} Runtime-S misses gold attributes; local gates cannot detect what was never extracted. The holdout shows nonzero external gold-literal exposure despite local PASS.
\item \textbf{Attacker scope.} Deterministic local audits and the external LLM audits do not cover arbitrary paraphrase, composition, or stronger attackers.
\item \textbf{Utility calibration.} The local utility surrogate is coarse; human and online-judge validation is limited and within-group correlations are modest.
\item \textbf{Neural local track not rerun.} The local LLM service was unavailable; new end-to-end results concern the deterministic local track.
\item \textbf{External audit disclosure.} External re-testing is controlled disclosure, not pre-egress prevention.
\item \textbf{Synthetic scope.} Results are on fictional synthetic benchmarks; real deployment requires independent data, stronger detectors, and domain-specific calibration.
\end{itemize}

\section{Conclusion}

DR-SL is best understood as an auditable cloud-local release workflow, not a certified zero-leakage sanitizer. The revised evidence supports deterministic truthful generalization plus all-occurrence scrubbing and stable local gating as useful components, while showing that runtime extraction misses create residual egress exposure. Theoretical statements are limited to leakage witnesses and MI-unit feasibility; statistical claims are scoped to observed attacks and instance clusters. Future work should prioritize stronger runtime extraction, local strong attackers or trusted execution environments, powered human calibration, and independent real-world anonymization benchmarks.

\paragraph{Artifact availability.} Code, synthetic datasets, frozen protocol, raw audit outputs, result hashes, and analysis scripts are public at \url{https://github.com/hurrywhen/dr-sl} and \url{https://gitee.com/jiangsu-yunhefeng-intelligent_0/dr-sl}. Dependencies are pinned in \texttt{requirements.txt}. API keys and private endpoints are excluded.

\appendix

\section*{Appendix A: Proof Sketches}

\subsection*{A.1 Fano-type bound}
Let $P^{*}_{e,i}=1-L_i^{*}$ be the MAP error. Fano's inequality gives $H(S_i\mid D)\le H_2(P^{*}_{e,i})+P^{*}_{e,i}\log(|V_i|-1)$. The function $g(p)=H_2(p)+p\log(k-1)$ is nondecreasing on $[0,(k-1)/k]$. Subadditivity $H(S\mid D)\le\sum_i H(S_i\mid D)$ yields the stated bound. No threshold substitution is valid: replacing a smaller actual hit rate by a larger threshold reverses the needed monotonicity direction.

\subsection*{A.2 Pufferfish witness}
Let $L^{*}$ be MAP success. There exists an output $d_0$ with posterior mode mass at least $L^{*}$. The posterior odds against any other value are at least $L^{*}/(1-L^{*})$. Bayes expansion relates posterior odds to likelihood ratio times prior odds, giving the stated lower bound on the mechanism's required $\varepsilon$.

\subsection*{A.3 Feasibility inequality}
Condition on $A$. Monotonicity and the chain rule give $I(D;T\mid A)\le I(D;S\mid A)+I(D;T\mid S,A)$. The conditional Markov relation $T\to X\to D$ given $(S,A)$ and the data-processing inequality give $I(D;T\mid S,A)\le I(X;T\mid S,A)$. The inequality is in mutual-information units only.

\subsection*{A.4 Termination}
Each transformation or skip increments a counter bounded by $B=3|S|+5$. Stability checks are bounded per artifact and artifact changes reset the window. The skip set is finite. Hence the algorithm terminates. Boundedness of a continuous utility score alone would not imply finitely many strict improvements.

\section*{Appendix B: Reproducibility}

The revised track is run by:
\begin{verbatim}
python3 dataset/generate_holdout.py
python3 sim/revision_experiment.py
python3 sim/revision_online_audit.py
python3 sim/make_revision_tables.py
python3 human_study/analyze_results.py
\end{verbatim}
Offline results require no LLM. Online audit calls use synthetic artifacts only and archive raw outputs in \texttt{results/revision\_20261118/online\_audit\_raw.jsonl}. The freeze manifest records SHA-256 hashes of protocol, generator, runtime code, and analysis scripts. Model snapshots and local weight quantization for the historical neural track are not fully pinned; this is a reproducibility limitation. HTTP backends now support an explicit seed field when configured, but historical runs did not submit seeds and cannot be claimed seed-reproducible.

\begin{thebibliography}{30}
\bibitem{hao2024hybrid} Hao Z, Jiang H, Jiang S, Ren J, Cao T. Hybrid SLM and LLM for edge-cloud collaborative inference. EdgeFM, 2024.
\bibitem{jin2025cecollm} Jin H, Wu Y. CE-COLLM: efficient and adaptive large language models through cloud-edge collaboration. IEEE ICWS, 2025.
\bibitem{li2025collab} Li S, Wang H, Xu W, et al. Collaborative inference and learning between edge SLMs and cloud LLMs: a survey. arXiv:2507.16731, 2025.
\bibitem{siyan2025papillon} Siyan L, Raghuram V C, Khattab O, et al. PAPILLON: privacy preservation from internet-based and local language model ensembles. NAACL, 2025.
\bibitem{das2025security} Das B C, Amini M H, Wu Y. Security and privacy challenges of large language models: a survey. ACM Computing Surveys, 2025.
\bibitem{lukas2023pii} Lukas N, Salem A, Sim R, et al. Analyzing leakage of personally identifiable information in language models. IEEE S\&P, 2023.
\bibitem{presidio} data-privacy-stack/presidio. Open-source framework for detecting, redacting, and anonymizing PII. GitHub.
\bibitem{rehydra} rehydra-ai/rehydra-sdk. PII security for AI workflows. GitHub, MIT.
\bibitem{zeng2025privacyrestore} Zeng Z, Wang J, Yang J, et al. PrivacyRestore: privacy-preserving inference in large language models via privacy removal and restoration. ACL, 2025.
\bibitem{mai2023split} Mai P, Yan R, Huang Z, Yang Y, Pang Y. Split-and-denoise: protect large language model inference with local differential privacy. arXiv:2310.09130, 2023.
\bibitem{shi2022selective} Shi W, Cui A, Li E, et al. Selective differential privacy for language modeling. NAACL, 2022.
\bibitem{duchi2013local} Duchi J C, Jordan M I, Wainwright M J. Local privacy and statistical minimax rates. IEEE FOCS, 2013.
\bibitem{yue2021santext} Yue X, Du M, Wang T, Li Y, Sun H, Chow S S M. Differential privacy for text analytics via natural text sanitization (SANTEXT). Findings of ACL/IJCNLP, 2021.
\bibitem{ryu2026p2skill} Ryu M, Park G, Lee S, Piao X, Kim J. P2Skill: privacy preserving skill distillation for cloud-local LLM inference systems. arXiv:2608.14094, 2026. v1: 2026-08-14; v2: 2026-10-13.
\bibitem{pilan2025truthful} Pil\'an I, Manzanares-Salor B, S\'anchez D, Lison P. Truthful text sanitization guided by inference attacks (INTACT). Applied Soft Computing, 185:114013, 2025. arXiv:2412.12928.
\bibitem{kifer2012rigorous} Kifer D, Machanavajjhala A. A rigorous and customizable framework for privacy. PODS, 2012.
\bibitem{kifer2014pufferfish} Kifer D, Machanavajjhala A. Pufferfish: a framework for mathematical privacy definitions. ACM TODS, 2014.
\bibitem{issa2020operational} Issa I, Wagner A B, Kamath S. An operational approach to information leakage. IEEE Trans. Inf. Theory, 2020.
\bibitem{sweeney2002k} Sweeney L. k-Anonymity: a model for protecting privacy. IJUFKS, 2002.
\bibitem{cover2006elements} Cover T M, Thomas J A. Elements of Information Theory. Wiley, 2006.
\bibitem{jin2021medqa} Jin Q, Dhingra B, Liu Z, Cohen W W, Lu X. What disease does this patient have? A large-scale open domain question answering dataset from medical exams. Applied Sciences, 2021.
\bibitem{wilson1927} Wilson E B. Probable inference, the law of succession, and statistical inference. JASA, 1927.
\end{thebibliography}

\end{document}
